%%=============================================================================
%% Methodologie
%%=============================================================================

\chapter{\IfLanguageName{dutch}{Methodologie}{Methodology}}%
\label{ch:methodologie}

%% TODO: In dit hoofstuk geef je een korte toelichting over hoe je te werk bent
%% gegaan. Verdeel je onderzoek in grote fasen, en licht in elke fase toe wat
%% de doelstelling was, welke deliverables daar uit gekomen zijn, en welke
%% onderzoeksmethoden je daarbij toegepast hebt. Verantwoord waarom je
%% op deze manier te werk gegaan bent.
%% 
%% Voorbeelden van zulke fasen zijn: literatuurstudie, opstellen van een
%% requirements-analyse, opstellen long-list (bij vergelijkende studie),
%% selectie van geschikte tools (bij vergelijkende studie, "short-list"),
%% opzetten testopstelling/PoC, uitvoeren testen en verzamelen
%% van resultaten, analyse van resultaten, ...
%%
%% !!!!! LET OP !!!!!
%%
%% Het is uitdrukkelijk NIET de bedoeling dat je het grootste deel van de corpus
%% van je bachelorproef in dit hoofstuk verwerkt! Dit hoofdstuk is eerder een
%% kort overzicht van je plan van aanpak.
%%
%% Maak voor elke fase (behalve het literatuuronderzoek) een NIEUW HOOFDSTUK aan
%% en geef het een gepaste titel.

In dit hoofdstuk wordt toegelicht hoe dit onderzoek wordt uitgevoerd. De aanpak bestaat uit verschillende fasen, waarbij zowel literatuurstudie als praktische data-analyse worden gecombineerd.
Het doel is om te komen tot een concreet stappenplan en een proof-of-concept dat aantoont hoe voetbaldata gestructureerd en bruikbaar kan worden ingezet binnen een clubcontext.

\section{Fase 1: Literatuurstudie}

In een eerste fase wordt relevante wetenschappelijke literatuur bestudeerd rond sport analytics, matchanalyse, spatio-temporele analyse en datagedreven besluitvorming in voetbal.
Deze literatuurstudie dient om:

\begin{itemize}
    \item inzicht te krijgen in de verschillende types voetbaldata (event data, tracking data, fysieke data),
    \item bestaande analysetechnieken en prestatie-indicatoren te begrijpen,
    \item typische implementatieproblemen binnen clubs in kaart te brengen.
\end{itemize}

Deze fase vormt de theoretische basis voor de verdere praktische uitwerking.

\section{Fase 2: Verkenning en structurering van de data}

In een tweede fase wordt de beschikbare data verkend.
Dit omvat in eerste instantie voorbeelddata (zoals publieke event- en trackingdata), en later de effectieve clubdata wanneer deze beschikbaar is.

In deze fase wordt:

\begin{itemize}
    \item de structuur van de datasets geanalyseerd (kolommen, variabelen, tijdsstructuur),
    \item nagegaan welke types informatie aanwezig zijn,
    \item gecontroleerd op ontbrekende waarden, inconsistenties of technische problemen.
\end{itemize}

Het doel van deze stap is om de data te begrijpen voordat er analyses worden uitgevoerd.

\section{Fase 3: Ontwikkeling van een proof-of-concept}

Op basis van de verkende data wordt een proof-of-concept (PoC) ontwikkeld in Python.
Deze PoC heeft als doel te tonen hoe ruwe voetbaldata kan worden:

\begin{itemize}
    \item opgeschoond en gestructureerd,
    \item omgezet naar bruikbare prestatie-indicatoren,
    \item gevisualiseerd op een begrijpelijke manier.
\end{itemize}

De focus ligt hierbij op haalbare en interpreteerbare analysetechnieken.
Voor event data kan dit bijvoorbeeld het berekenen van kanskwaliteit of actie-evaluatie zijn.
Voor tracking data kan dit het analyseren van ruimtegebruik of veldbezetting omvatten.

Het doel is niet om een complex model te bouwen, maar om een reproduceerbare en praktische analysemethode te demonstreren.

\section{Fase 4: Analyse van implementatie binnen clubcontext}

Naast de technische analyse wordt ook onderzocht hoe dergelijke analysemethoden praktisch kunnen worden ingebed binnen een voetbalclub.
Hierbij wordt gekeken naar:

\begin{itemize}
    \item benodigde tijd en expertise,
    \item technische vereisten,
    \item mogelijke knelpunten bij integratie in bestaande workflows.
\end{itemize}

Deze reflectie is belangrijk om ervoor te zorgen dat het eindresultaat niet louter theoretisch blijft.

\section{Fase 5: Uitwerking van een stappenplan}

Op basis van de voorgaande fasen wordt een concreet stappenplan opgesteld.
Dit stappenplan beschrijft hoe een voetbalclub gefaseerd kan evolueren naar een meer gestructureerde en efficiënte inzet van data.

Het stappenplan omvat onder meer:

\begin{itemize}
    \item het organiseren en centraliseren van databronnen,
    \item het definiëren van duidelijke analysevragen,
    \item het kiezen van geschikte tools,
    \item het vertalen van analyses naar bruikbare rapportering.
\end{itemize}

Het einddoel van deze bachelorproef is dus tweeledig: enerzijds het aantonen van een werkbare analysemethode via een proof-of-concept, anderzijds het formuleren van praktische aanbevelingen voor implementatie binnen een voetbalclub.